\item Encuentre ecuaciones paramétricas para describir el conjunto de todos los puntos $P$ determinados como se ilustra en la figura, utilizando el ángulo $\theta$ como parámetro. El segmento de recta $AB$ es tangente al círculo mayor.
    
    \begin{center}
    \begin{tikzpicture}[scale=1.2]
        \draw[->] (-1.5,0) -- (3.5,0) node[below] {$x$};
        \draw[->] (0,-1.5) -- (0,2.5) node[left] {$y$};
        
        \draw[thick] (0,0) circle (1cm); 
        \draw[thick] (0,0) circle (2cm); 
        \node[below left] at (0,0) {$0$};
        
        \coordinate (A) at (50:2);
        \draw (0,0) -- (A);
        
        \coordinate (B) at ({2/cos(50)}, 0);
        \draw[thick] (A) -- (B) node[below] {$B$};
        \node[above left] at (A) {$A$};
        
        \coordinate (P) at ({2/cos(50)}, {2*sin(50)});
        \draw[thick] (A) -- (P);
        \draw[thick] (B) -- (P);
        
        \fill (P) circle (1pt) node[above right] {$P$};
        
        \node[above left] at (140:2) {$a$};
        \node[above left] at (140:1) {$b$};
        
        \coordinate (origin) at (0,0);
        \pic [draw, "$\theta$", angle eccentricity=1.5, angle radius=0.5cm] {angle = B--origin--A};
    \end{tikzpicture}
    \end{center}
